\pagestyle{empty}
\tight
\begin{tblr}{width=\textwidth,colspec={|Q[c,m,1.9cm]|X[l,m]|},hlines,vlines,hspan=minimal,row{1}={bg=headblue},rowsep=1pt,colsep=3pt}
\SetCell{c,m}\hfil\logo\hfil & \textbf{POLITEKNIK NEGERI MALANG}\par\textbf{JURUSAN TEKNOLOGI INFORMASI}\par\textbf{PROGRAM STUDI: D4 TEKNIK INFORMATIKA} \\
\SetCell[c=2]{c} \textbf{RENCANA PEMBELAJARAN SEMESTER (RPS)} & \\
\end{tblr}
\begin{longtblr}{width=\textwidth,colspec={|Q[l,m,1.9cm]|X[0.85,l,m]|X[1.45,l,m]|X[0.75,l,m]|X[1.15,l,m]|},hlines,vlines,hspan=minimal,rowsep=1pt,colsep=2pt,row{1,3}={bg=headgray,font=\bfseries}}
MATA KULIAH & KODE & BOBOT (SKS)/jam & SEMESTER & TGL. PENYUSUNAN \\
Komunikasi dan Etika Profesi & RTI256001 & 2 SKS/4 jam & 6 & 1 Juli 2025 \\
OTORISASI & Dosen Pengembang RPS & Koordinator RMK & \SetCell[c=2]{l} Ka PRODI &  \\
 & Triana Fatmawati, S.T., M.T. & Triana Fatmawati, S.T., M.T. & \SetCell[c=2]{l}{\vspace{8mm}\par Dr. Ely Setyo Astuti, S.T., M.T.} &  \\
\SetCell[r=11]{l} Capaian Pembelajaran (CP) & \SetCell[c=4]{l,bg=headgray,font=\bfseries} Capaian Pembelajaran Lulusan Program Studi (CPL-Prodi) &  &  &  \\
 & CPL01 & \SetCell[c=3]{l} Mampu menerapkan nilai ketakwaan, kesadaran hukum, kedisiplinan, profesionalisme, etika profesi, dan pembelajaran sepanjang hayat, serta tanggap terhadap isu sosial dan perkembangan teknologi. &  &  \\
 & CPL03 & \SetCell[c=3]{l} Mampu mengelola tim dan diri sendiri secara profesional serta mengomunikasikan dan mempresentasikan gagasan maupun hasil kerja secara efektif baik lisan maupun tulisan. &  &  \\
 & \SetCell[c=4]{l,bg=headgray,font=\bfseries} Capaian Pembelajaran mata kuliah (CPL-MK) &  &  &  \\
 & CPMK0101 & \SetCell[c=3]{l} Mahasiswa mampu menerapkan etika profesi TI, kesadaran hukum, K3, dan nilai profesionalisme dalam lingkungan kerja digital. &  &  \\
 & CPMK0301 & \SetCell[c=3]{l} Mahasiswa mampu berkomunikasi dan mempresentasikan gagasan teknis secara efektif dalam konteks kerja profesional. &  &  \\
 & CPMK0303 & \SetCell[c=3]{l} Mahasiswa mampu menyusun rencana pengembangan karir dan portofolio profesional di bidang TI. &  &  \\
 & \SetCell[c=4]{l,bg=headgray,font=\bfseries} Kemampuan akhir tiap tahapan belajar (Sub-CPMK) &  &  &  \\
 & SCPMK0101-04501 & \SetCell[c=3]{l} Mahasiswa mampu menganalisis kasus pelanggaran etika profesi, hukum ITE, dan K3 untuk merekomendasikan tindakan yang tepat. &  &  \\
 & SCPMK0301-04502 & \SetCell[c=3]{l} Mahasiswa mampu menyajikan presentasi teknis, menulis laporan profesional, dan berkolaborasi dalam tim lintas fungsi. &  &  \\
 & SCPMK0303-04503 & \SetCell[c=3]{l} Mahasiswa mampu menyusun portofolio profesional dan rencana pengembangan karir berbasis kompetensi TI. &  &  \\
\SetCell[r=2]{l} Penilaian & \SetCell[c=4]{l} *Nilai CPMK didapat dari agregasi nilai SUB-CPMK yang dibebankan &  &  &  \\
 & \SetCell[c=4]{l}{\tiny\begin{tblr}{width=\linewidth,colspec={|Q[c,m,0.1\linewidth]|Q[c,m,0.16\linewidth]|X[l,m]|X[l,m]|X[l,m]|X[l,m]|X[l,m]|X[l,m]|Q[c,m,0.08\linewidth]|},hlines,vlines,hspan=minimal,rowsep=1pt,colsep=0.7pt,row{1,2}={bg=headgray,font=\bfseries}}
Kode CPMK & Kode SCPMK & \SetCell[c=6]{c} Bobot per Bentuk Penilaian & & & & & & Total Bobot \\
 & & Tugas & Kuis 1 & UTS & Kuis 2 & PBL & UAS & \\
CPMK0101 & SCPMK0101-04501 & 14\%: etika dan hukum ITE & 5\%: etika dan hukum & 7\%: kasus etika & 0\% & 0\% & 0\% & 26\% \\
CPMK0301 & SCPMK0301-04502 & 10\%: komunikasi teknis & 0\% & 6\%: komunikasi profesional & 0\% & 10\%: simulasi profesional & 0\% & 26\% \\
CPMK0303 & SCPMK0303-04503 & 6\%: perencanaan karir & 0\% & 0\% & 5\%: karir TI & 12\%: portofolio & 25\%: presentasi portofolio & 48\% \\
\SetCell[c=8]{r}\textbf{Total Per Penilaian} & & & & & & & & \textbf{100\%} \\
\end{tblr}} &  &  &  \\
Diskripsi Singkat Mata Kuliah & \SetCell[c=4]{l} Mata kuliah Komunikasi dan Etika Profesi membangun kompetensi etika profesi TI, komunikasi profesional, dan perencanaan karir melalui studi kasus, simulasi wawancara, dan penyusunan portofolio digital yang siap kerja. &  &  &  \\
Bahan Kajian / Materi Pembelajaran & \SetCell[c=4]{l}\begin{enumerate}\item Etika profesi TI, kode etik ACM/IEEE, hukum ITE, dan K3\item Komunikasi profesional: presentasi teknis dan penulisan laporan\item Dinamika tim, kolaborasi lintas fungsi, dan isu sosial teknologi\item Personal branding, portofolio digital, dan perencanaan karir\item Strategi pencarian kerja, CV, wawancara, dan etika penelitian\end{enumerate} &  &  &  \\
\SetCell[r=2]{l} Pustaka & \SetCell[c=4]{l} \textbf{Utama:}\begin{enumerate}\item ACM/IEEE Code of Ethics 2018.\item Locker, K. 2020. Business Communication. McGraw-Hill.\item Floridi, L. 2014. The Ethics of AI. MIT Press.\end{enumerate} &  &  &  \\
 & \SetCell[c=4]{l} \textbf{Pendukung:} Materi LMS, dokumentasi resmi, dan jobsheet praktikum. &  &  &  \\
Media Pembelajaran & \SetCell[c=2]{l} \textbf{Software:} Microsoft Office, LMS, LinkedIn, platform portofolio digital. &  & \SetCell[c=2]{l} \textbf{Hardware:} Komputer/laptop, proyektor. &  \\
Nama Dosen Pengampu & \SetCell[c=4]{l} Tim Pengajar Program Studi D4 TI &  &  &  \\
Matakuliah Syarat & \SetCell[c=4]{l} - &  &  &  \\
\end{longtblr}
\begin{longtblr}{width=\textwidth,colspec={|Q[c,m,0.06\textwidth]|X[1.35,l,m]|X[1.08,l,m]|X[1.16,l,m]|X[0.7,l,m]|X[1.24,l,m]|X[1.06,l,m]|X[0.98,l,m]|Q[c,m,0.05\textwidth]|},hlines,vlines,hspan=minimal,rowhead=2,row{1,2}={bg=headgreen,font=\bfseries},rowsep=1pt,colsep=1.5pt}
Mgg.\par Ke & Kemampuan Akhir Yang Direncanakan (Sub-CP-MK) & Bahan kajian (Materi Pembelajaran) & Bentuk dan Metode Pembelajaran & Estimasi Waktu & Pengalaman Belajar Mahasiswa & Kriteria \& Bentuk Penilaian & Indikator Penilaian & Bobot Penilaian (\%) \\
(1) & (2) & (3) & (4) & (5) & (6) & (7) & (8) & (9) \\
1 & SCPMK0101-04501 & Pengantar etika profesi TI dan kode etik ACM/IEEE. & Modalitas: Blended Learning. Bentuk: Luring/Daring. Metode: Case Method, simulasi profesional, studi kasus, dan diskusi reflektif. & 1 x 2 x 50' tatap muka; 1 x 2 x 50' tugas/praktik mandiri. & Mahasiswa mengerjakan aktivitas terarah untuk pengantar etika profesi TI dan kode etik ACM/IEEE dan menunjukkan evidence hasil belajar. & Bentuk penilaian: Pengantar Etika Profesi TI dan Kode Etik ACM/IEEE. Mengacu rubrik RTM. & Ketepatan konsep; kualitas komunikasi; validasi dan dokumentasi. & 2 \\
2 & SCPMK0101-04501 & Hukum ITE dan regulasi perlindungan data digital. & Modalitas: Blended Learning. Bentuk: Luring/Daring. Metode: Case Method, simulasi profesional, studi kasus, dan diskusi reflektif. & 1 x 2 x 50' tatap muka; 1 x 2 x 50' tugas/praktik mandiri. & Mahasiswa mengerjakan aktivitas terarah untuk hukum ITE dan regulasi perlindungan data digital dan menunjukkan evidence hasil belajar. & Bentuk penilaian: Hukum ITE dan Regulasi Perlindungan Data Digital. Mengacu rubrik RTM. & Ketepatan konsep; kualitas komunikasi; validasi dan dokumentasi. & 3 \\
3 & SCPMK0101-04501 & K3 di lingkungan kerja TI: ergonomi dan keselamatan digital. & Modalitas: Blended Learning. Bentuk: Luring/Daring. Metode: Case Method, simulasi profesional, studi kasus, dan diskusi reflektif. & 1 x 2 x 50' tatap muka; 1 x 2 x 50' tugas/praktik mandiri. & Mahasiswa mengerjakan aktivitas terarah untuk K3 di lingkungan kerja TI: ergonomi dan keselamatan digital dan menunjukkan evidence hasil belajar. & Bentuk penilaian: K3 di Lingkungan Kerja TI: Ergonomi dan Keselamatan Digital. Mengacu rubrik RTM. & Ketepatan konsep; kualitas komunikasi; validasi dan dokumentasi. & 3 \\
4 & SCPMK0101-04501 & Kuis 1: etika dan hukum TI. & Modalitas: Blended Learning. Bentuk: Luring/Daring. Metode: Case Method, simulasi profesional, studi kasus, dan diskusi reflektif. & 1 x 2 x 50' tatap muka; 1 x 2 x 50' tugas/praktik mandiri. & Mahasiswa mengerjakan aktivitas terarah untuk kuis 1: etika dan hukum TI dan menunjukkan evidence hasil belajar. & Bentuk penilaian: Kuis 1: Etika dan Hukum TI. Mengacu rubrik RTM. & Ketepatan konsep; kualitas komunikasi; validasi dan dokumentasi. & 5 \\
5 & SCPMK0301-04502 & Teknik komunikasi profesional dan presentasi teknis. & Modalitas: Blended Learning. Bentuk: Luring/Daring. Metode: Case Method, simulasi profesional, studi kasus, dan diskusi reflektif. & 1 x 2 x 50' tatap muka; 1 x 2 x 50' tugas/praktik mandiri. & Mahasiswa mengerjakan aktivitas terarah untuk teknik komunikasi profesional dan presentasi teknis dan menunjukkan evidence hasil belajar. & Bentuk penilaian: Teknik Komunikasi Profesional dan Presentasi Teknis. Mengacu rubrik RTM. & Ketepatan konsep; kualitas komunikasi; validasi dan dokumentasi. & 3 \\
6 & SCPMK0301-04502 & Penulisan profesional: email, laporan, dan proposal teknis. & Modalitas: Blended Learning. Bentuk: Luring/Daring. Metode: Case Method, simulasi profesional, studi kasus, dan diskusi reflektif. & 1 x 2 x 50' tatap muka; 1 x 2 x 50' tugas/praktik mandiri. & Mahasiswa mengerjakan aktivitas terarah untuk penulisan profesional: email, laporan, dan proposal teknis dan menunjukkan evidence hasil belajar. & Bentuk penilaian: Penulisan Profesional: Email, Laporan, dan Proposal Teknis. Mengacu rubrik RTM. & Ketepatan konsep; kualitas komunikasi; validasi dan dokumentasi. & 3 \\
7 & SCPMK0301-04502 & Dinamika tim dan kolaborasi lintas fungsi. & Modalitas: Blended Learning. Bentuk: Luring/Daring. Metode: Case Method, simulasi profesional, studi kasus, dan diskusi reflektif. & 1 x 2 x 50' tatap muka; 1 x 2 x 50' tugas/praktik mandiri. & Mahasiswa mengerjakan aktivitas terarah untuk dinamika tim dan kolaborasi lintas fungsi dan menunjukkan evidence hasil belajar. & Bentuk penilaian: Dinamika Tim dan Kolaborasi Lintas Fungsi. Mengacu rubrik RTM. & Ketepatan konsep; kualitas komunikasi; validasi dan dokumentasi. & 3 \\
8 & SCPMK0101-04501, SCPMK0301-04502 & UTS: kasus etika, hukum, dan komunikasi. & Modalitas: Blended Learning. Bentuk: Luring/Daring. Metode: Case Method, simulasi profesional, studi kasus, dan diskusi reflektif. & 1 x 2 x 50' tatap muka; 1 x 2 x 50' tugas/praktik mandiri. & Mahasiswa mengerjakan aktivitas terarah untuk UTS: kasus etika, hukum, dan komunikasi dan menunjukkan evidence hasil belajar. & Bentuk penilaian: UTS: Kasus Etika, Hukum, dan Komunikasi. Mengacu rubrik RTM. & Ketepatan konsep; kualitas komunikasi; validasi dan dokumentasi. & 13 \\
9 & SCPMK0101-04501 & Isu sosial teknologi: privasi data dan surveillance. & Modalitas: Blended Learning. Bentuk: Luring/Daring. Metode: Case Method, simulasi profesional, studi kasus, dan diskusi reflektif. & 1 x 2 x 50' tatap muka; 1 x 2 x 50' tugas/praktik mandiri. & Mahasiswa mengerjakan aktivitas terarah untuk isu sosial teknologi: privasi data dan surveillance dan menunjukkan evidence hasil belajar. & Bentuk penilaian: Isu Sosial Teknologi: Privasi Data dan Surveillance. Mengacu rubrik RTM. & Ketepatan konsep; kualitas komunikasi; validasi dan dokumentasi. & 2 \\
10 & SCPMK0303-04503 & Tanggung jawab sosial korporasi di bidang TI. & Modalitas: Blended Learning. Bentuk: Luring/Daring. Metode: Case Method, simulasi profesional, studi kasus, dan diskusi reflektif. & 1 x 2 x 50' tatap muka; 1 x 2 x 50' tugas/praktik mandiri. & Mahasiswa mengerjakan aktivitas terarah untuk tanggung jawab sosial korporasi di bidang TI dan menunjukkan evidence hasil belajar. & Bentuk penilaian: Tanggung Jawab Sosial Korporasi di Bidang TI. Mengacu rubrik RTM. & Ketepatan konsep; kualitas komunikasi; validasi dan dokumentasi. & 2 \\
11 & SCPMK0303-04503 & Personal branding: LinkedIn, GitHub, dan portofolio digital. & Modalitas: Blended Learning. Bentuk: Luring/Daring. Metode: Case Method, simulasi profesional, studi kasus, dan diskusi reflektif. & 1 x 2 x 50' tatap muka; 1 x 2 x 50' tugas/praktik mandiri. & Mahasiswa mengerjakan aktivitas terarah untuk personal branding: LinkedIn, GitHub, dan portofolio digital dan menunjukkan evidence hasil belajar. & Bentuk penilaian: Personal Branding: LinkedIn, GitHub, dan Portofolio Digital. Mengacu rubrik RTM. & Ketepatan konsep; kualitas komunikasi; validasi dan dokumentasi. & 3 \\
12 & SCPMK0303-04503 & Kuis 2: isu sosial dan karir TI. & Modalitas: Blended Learning. Bentuk: Luring/Daring. Metode: Case Method, simulasi profesional, studi kasus, dan diskusi reflektif. & 1 x 2 x 50' tatap muka; 1 x 2 x 50' tugas/praktik mandiri. & Mahasiswa mengerjakan aktivitas terarah untuk kuis 2: isu sosial dan karir TI dan menunjukkan evidence hasil belajar. & Bentuk penilaian: Kuis 2: Isu Sosial dan Karir TI. Mengacu rubrik RTM. & Ketepatan konsep; kualitas komunikasi; validasi dan dokumentasi. & 5 \\
13 & SCPMK0303-04503 & Perencanaan karir dan strategi pencarian kerja. & Modalitas: Blended Learning. Bentuk: Luring/Daring. Metode: Case Method, simulasi profesional, studi kasus, dan diskusi reflektif. & 1 x 2 x 50' tatap muka; 1 x 2 x 50' tugas/praktik mandiri. & Mahasiswa mengerjakan aktivitas terarah untuk perencanaan karir dan strategi pencarian kerja dan menunjukkan evidence hasil belajar. & Bentuk penilaian: Perencanaan Karir dan Strategi Pencarian Kerja. Mengacu rubrik RTM. & Ketepatan konsep; kualitas komunikasi; validasi dan dokumentasi. & 2 \\
14 & SCPMK0303-04503 & CV, cover letter, dan teknik wawancara kerja. & Modalitas: Blended Learning. Bentuk: Luring/Daring. Metode: Case Method, simulasi profesional, studi kasus, dan diskusi reflektif. & 1 x 2 x 50' tatap muka; 1 x 2 x 50' tugas/praktik mandiri. & Mahasiswa mengerjakan aktivitas terarah untuk CV, cover letter, dan teknik wawancara kerja dan menunjukkan evidence hasil belajar. & Bentuk penilaian: CV, Cover Letter, dan Teknik Wawancara Kerja. Mengacu rubrik RTM. & Ketepatan konsep; kualitas komunikasi; validasi dan dokumentasi. & 2 \\
15 & SCPMK0101-04501 & Etika penelitian dan publikasi ilmiah. & Modalitas: Blended Learning. Bentuk: Luring/Daring. Metode: Case Method, simulasi profesional, studi kasus, dan diskusi reflektif. & 1 x 2 x 50' tatap muka; 1 x 2 x 50' tugas/praktik mandiri. & Mahasiswa mengerjakan aktivitas terarah untuk etika penelitian dan publikasi ilmiah dan menunjukkan evidence hasil belajar. & Bentuk penilaian: Etika Penelitian dan Publikasi Ilmiah. Mengacu rubrik RTM. & Ketepatan konsep; kualitas komunikasi; validasi dan dokumentasi. & 2 \\
16 & SCPMK0301-04502, SCPMK0303-04503 & PBL: simulasi wawancara dan presentasi portofolio. & Modalitas: Blended Learning. Bentuk: Luring/Daring. Metode: Case Method, simulasi profesional, studi kasus, dan diskusi reflektif. & 1 x 2 x 50' tatap muka; 1 x 2 x 50' tugas/praktik mandiri. & Mahasiswa mengerjakan aktivitas terarah untuk PBL: simulasi wawancara dan presentasi portofolio dan menunjukkan evidence hasil belajar. & Bentuk penilaian: PBL: Simulasi Wawancara dan Presentasi Portofolio. Mengacu rubrik RTM. & Ketepatan konsep; kualitas komunikasi; validasi dan dokumentasi. & 22 \\
17 & SCPMK0303-04503 & UAS: presentasi portofolio dan rencana karir. & Modalitas: Blended Learning. Bentuk: Luring/Daring. Metode: Case Method, simulasi profesional, studi kasus, dan diskusi reflektif. & 1 x 2 x 50' tatap muka; 1 x 2 x 50' tugas/praktik mandiri. & Mahasiswa mengerjakan aktivitas terarah untuk UAS: presentasi portofolio dan rencana karir dan menunjukkan evidence hasil belajar. & Bentuk penilaian: UAS: Presentasi Portofolio dan Rencana Karir. Mengacu rubrik RTM. & Ketepatan konsep; kualitas komunikasi; validasi dan dokumentasi. & 25 \\
\end{longtblr}
